% TonalWM — ICASSP 2027 submission (4 pages + references)
% ==========================================================================
% This is the CUT-DOWN version. The complete study lives in paper/ojsp_full.tex
% and goes to OJSP (SPEC §9 plan, unchanged). Material cut to fit — ambiguity
% strata, specificity matrix, fifths-distance curve, equivariance per-layer,
% surgical-control curve — is referenced as supplementary, never deleted.
%
% INTEGRITY (SPEC §7): every number traces to a ledgered artifact under results/.
% No number is estimated. "TBD (run required)" marks an unmade measurement.
%
% FORMAT: written for spconf.sty (ICASSP Paper Kit, not redistributable); falls
% back to IEEEtran so the draft always compiles. No content depends on which.
% ==========================================================================
\documentclass[conference]{IEEEtran}
\newif\ifspconf
\IfFileExists{spconf.sty}{\spconftrue}{\spconffalse}
\ifspconf \usepackage{spconf} \fi
\ifspconf\else
  \newcommand{\name}[1]{\author{\IEEEauthorblockN{#1}}}
  \newcommand{\address}[1]{}
  \newcommand{\ninept}{}
  \newenvironment{keywords}{\begin{IEEEkeywords}}{\end{IEEEkeywords}}
\fi
\usepackage[T1]{fontenc}
\usepackage{lmodern}
\usepackage{amsmath,amssymb,graphicx,booktabs,array,tabularx}
\usepackage[hidelinks]{hyperref}
\usepackage{url}
\graphicspath{{../results/figures/}}

\newcommand{\tkr}{\mathrm{TKR}}
\newcommand{\ikr}{\mathrm{IKR}}
\newcommand{\ecyc}{\varepsilon_{\mathrm{cyc}}}
\newcommand{\ppl}{\delta_{\mathrm{PPL}}}
\newcommand{\TBD}[1]{\textbf{[TBD (run required): #1]}}
\newcommand{\vs}{vs.\@\ }

\title{Does a Music Transformer Know What Key It Is In?\\
Causal Interventions on an Emergent Tonal World Model}

\name{Anonymous ICASSP 2027 submission}
\address{Paper ID: TBD}

\begin{document}
\ninept
\maketitle

\begin{abstract}
Music generation models produce convincingly tonal music. Whether they chain
surface statistics or maintain an internal representation of musical key that
generation \emph{causally uses} has not been tested. We test it directly, under a
pre-registered protocol with frozen decision rules, on Transformers trained on
symbolic music whose vocabulary holds absolute pitch and time only---no key,
chord, or degree tokens. Replacing the key component of the residual stream at a
single middle layer makes the model continue in the injected key at $5.0\times$
the rate of rank- and norm-matched random edits (guarded $\tkr$ $0.378$ \vs
$0.075$; significant for $12/12$ targets after Holm correction, replicated across
seeds), inside a musicality budget frozen before any edit ran, reaching $58\%$ of
the ceiling set by transposing the prompt itself. Two controls say what this
requires. Across a $170\times$ capacity range next-token accuracy is flat, yet at
$0.5$M parameters the key is not decodable beyond the input surface at all and the
same edit moves the key \emph{more} often than at $26$M while wrecking the music
$90\%$ of the time: capacity does not buy the key, it buys a key that can be moved
without breaking everything else. Finally we leave our own models: on real Bach chorales a
\emph{public} checkpoint trained on real music---architecturally identical to one of
ours---carries a key state that beats the pitch surface where ours does not, and editing
that state redirects the music it composes at $5.7\times$ the matched control ($11/12$
targets, $100\%$ within the frozen musicality budget). It also resists: the keys the edit
cannot reach are the keys real music rarely uses. Editing a music Transformer's internal
key state changes what key it composes in---in a model people actually use, not only in
ours.
\end{abstract}

\begin{keywords}
music generation, world models, interpretability, causal intervention, tonality
\end{keywords}

\section{Introduction}
\label{sec:intro}
Systems that compose music write convincingly tonal pieces: they settle into a
home key, build tension away from it, and return. How they do so is unexamined.
Analyses report that musical attributes can be \emph{read out} of a model's
activations by trained classifiers~\cite{syntheory2024}. Reading is not using: a
signal can be decodable from a system's state without the system consulting it
when it acts~\cite{hewitt2019control}. No study has tested whether a music model
\emph{uses} an internal key representation---whether intervening on it, with the
input untouched, changes the key of the music it goes on to generate.

Sequence models can build causal models of the latent state behind their data. A
model trained only to predict Othello moves tracks the board those moves imply,
and editing that internal board changes which moves it treats as
legal~\cite{li2023othello,nanda2023emergent}; analogous causal evidence exists for
chess, mazes, and RL environments~\cite{karvonen2024chess,maze2024,iris2026}. In
music, evidence has stopped at readability~\cite{syntheory2024}, and steering work
controls attributes without testing a state
variable~\cite{smitin2024}. We occupy the empty cell: causal, state-level evidence
in music.

Key is an unusually good state variable for this. It is defined \emph{externally}
to the model---by centuries of theory and by psychological
measurement~\cite{krumhansl1982}---so the probed state is not an artifact of our
labeling; it must be \emph{computed} by integrating pitch evidence over time,
unlike metric position, which time tokens hand over; and diatonic membership gives
a natural analogue of Othello's legal-move rate, so success is measurable without
a learned judge.

Contributions: (1)~a controlled testbed---leak-free tokenizer, functional-harmony
corpus with exact per-token key labels, pre-registered protocol---released as a
benchmark; (2)~a causal result: single-layer subspace edits redirect the composed
key for $12/12$ targets inside a frozen musicality budget; (3)~two controls that
say what a world model requires: a capacity sweep separating the state's presence,
its legibility, and its manipulability; and the same probe AND the same intervention
carried to a public model trained on real music, which reads its key state and uses it.

\section{Method}
\label{sec:method}

\subsection{Data, models, protocol}
Ground truth must be exact, not estimated, so training data are synthetic.
\textbf{D-SYN}: a functional-harmony state machine (T$\to$S$\to$D$\to$T with
substitutions, SATB voicing, optional diatonic melody) generates
$200$k$/10$k$/10$k train/val/test pieces of $278$--$508$ tokens with $0$--$3$
modulations each (pivot, direct, or sequential; fifths-distance stratified), giving
every token a ground-truth key label $\kappa(t)$ (tonic $\times$ mode, $24$
classes). \textbf{Tokenizer}: the vocabulary's $124$ types are \texttt{BAR},
\texttt{POS}, absolute-MIDI \texttt{PITCH} and \texttt{DUR} ($121$ types) together with
\texttt{PAD}/\texttt{BOS}/\texttt{EOS}; nothing names a key, chord or degree, and a unit
test enforces this, so key must be computed from the notes.
\textbf{Models}: M-CTRL is a GPT-2 style decoder ($8$ layers, $8$ heads,
$d{=}512$, ctx $512$, ${\sim}25$M params), trained under two regimes---with
(R-Aug) and without (R-NoAug) uniform $12$-fold transposition
augmentation---$\times\,3$ seeds. M-REF shares the architecture but a disjoint seed
\emph{and} data split, and serves only as the musicality guard, avoiding circular
evaluation. All models pass a pre-frozen quality gate (val.\ next-token top-$1$
$0.878$--$0.881$ \vs a $0.434$ constant-predictor rate). Decision rules, thresholds
and controls were frozen before any experiment ran; deviations are logged and
negative results reported.

\subsection{Probing (Phase A)}
Linear probes map residual-stream activations $h_\ell(t)$ to $\kappa(t)$ with
sequence-level splits, against three control families: \textbf{C1}, a control
task~\cite{hewitt2019control} permuting key identities per sequence; \textbf{C2},
an untrained model; \textbf{C3}, input-side baselines at the same positions
(pitch-class-histogram logistic regression and Krumhansl--Schmuckler (KS) template
matching, windows $W\in\{8,16,32,64\}$). \textbf{DR-H1}: the probe supports H1 iff
its selectivity-corrected $F_1$ exceeds the best C3 baseline with a sequence-level
BCa $95\%$ CI excluding zero in some layer.

The pre-registered hypotheses are H1 (key is decodable above the input surface), H2
(the representation is transposition-equivariant with cyclic-group structure) with
sub-claim H2b (augmentation strengthens equivariance---hence no ``H2a'': H2 is the
parent), H3 (a targeted subspace edit moves the composed key, beating a rank- and
norm-matched random edit inside the guard), H4 (after a \emph{one-shot} edit the state
persists, is reasserted, or transitions via a pivot), and H5 (causally effective layers
are localized). H4 is Phase~C: the pre-registration assigned it to the journal version
before any run, so it is out of scope here by design rather than for space, and it has
not been run.

\subsection{Intervention (Phase B)}
\label{sec:methodB}
From bar~$9$ ($t^{*}$) of a key-stable $8$-bar prompt, a sustained edit replaces the key
component at one layer at \emph{every} token position from $t^{*}$ on (not only where a
pitch is chosen---we write to the whole stream but probe one position in it), at every
generation step,
\begin{equation}
h \;\leftarrow\; h - P_V h + P_V\,\mu_{\kappa^{*}},
\label{eq:edit}
\end{equation}
with $V$ an orthonormal key subspace, $P_V=VV^{\!\top}$, and $\mu_{\kappa^{*}}$ the
class-conditional mean of the donor key $\kappa^{*}$. Subspaces: \textbf{V-PROBE}
(probe row space, rank $24$), \textbf{V-MEAN} (class-mean span), \textbf{V-DAS}
(interchange-trained, ranks $8/24$). Controls: \textbf{K1} a random subspace matched in
rank and norm---the bases are orthonormal, so $\lVert V\rVert_F\!=\!\sqrt{24}$ for both,
and the perturbations they apply agree to under $1\%$ when measured on real activations
($\lVert\Delta\rVert$ $21.9$ \vs $22.1$, the control marginally the larger), so what
separates them is direction, not magnitude; \textbf{K2} sham edit, required bit-identical to the clean run (passed,
$100/100$); \textbf{K3} another layer's subspace; \textbf{K4} the behavioral
ceiling---transposing the prompt itself. Continuations ($16$ bars; $12$ major
targets $\times\,100$ prompts $\times\,8$ layers; sampling noise paired with a clean
twin) are scored by KS key estimation ($\tkr$ strict: the continuation's estimated key
equals $\kappa^{*}$; all headline numbers are strict) and by $\ikr$ (fraction of
continuation pitches diatonic to a key). The pre-registered \emph{tolerant} variant,
used only to bracket estimator confusions, counts fifths distance ${\leq}1$ with the
same mode, the relative, or the parallel key. The musicality budget $\ppl{=}0.613$---the $90$th percentile of the M-REF
perplexity rise across $7{,}989$ \emph{natural} modulations in validation
data---was fixed \emph{before} any edit ran; edits exceeding it do not count as
successes. \textbf{DR-H3} requires $\tkr(\text{edit})>\tkr(\text{K1})$ with
Holm-corrected Wilcoxon $p<.05$ in ${\geq}\,8/12$ targets at the best in-budget
condition.

\begin{figure}[t]
\centering
\includegraphics[width=0.9\columnwidth]{fig_framework.pdf}
\caption{The claim in one example (real model output, fixed sampling seed). The same
$8$-bar B$\flat$-major prompt, continued (a)~clean and (b)~with the key component of
layer-$4$ activations replaced by the E-major class mean from bar~$9$ on.
\textbf{The input does not change}; the continuation's estimated key follows the
injected internal state. Shading marks the scale of the key in force.}
\label{fig:overview}
\end{figure}

\section{Results}
\label{sec:results}
Intervention numbers are from the R-Aug seed-$0$ sweep; probing holds for all six
M-CTRL models, and the seed-$1$ sweep replicates every Phase-B verdict, though seed~$1$
peaks at L$2$ ($\tkr$ $0.321$ \vs K1 $0.056$): the localization replicates, its depth
does not.

\subsection{The key is decodable beyond the input surface}
\label{sec:h1}
The key is \emph{computed}, not transcribed: at L$4$ the linear probe reaches
macro-$F_1$ $0.927$ against $0.790$ for the best input baseline (PC-hist LR,
$W{=}64$) and a $0.032$ selectivity floor, so the corrected margin is $+0.105$ (BCa
$95\%$ CI $[0.091,0.114]$). The margin excludes zero in layers $1$--$7$ of every
trained model. It fails only at L$0$, where the margin is significantly \emph{negative}
($-0.20$, CI $[-0.211,-0.182]$): at the embedding the notes themselves beat the
network's own read of them, exactly as they should. DR-H1 is supported $6/6$. An
untrained model scores $0.243$. The internal state is worth most where the surface
is worst: in the top quartile of KS-entropy ambiguity, probe accuracy is $0.936$
\vs $0.795$ for the baseline.

\subsection{Fifths geometry; equivariant but not a group}
The representation carries music-theoretic structure. Projected to two principal
components, the twelve major-tonic probe directions order themselves as
C--G--D--A--E--B--F$\sharp$--$\cdots$ (Fig.~\ref{fig:fifths}): circle-of-fifths
neighbour agreement $1.00$ \vs $0.00$ for the chromatic ordering, at every probed layer
(exploratory). Transposition acts as a generalizing linear map---Procrustes $R_k$ fit on
half the sequences aligns held-out activations with $0.05$--$0.11$ relative error---but
the family is not cyclic ($\ecyc\approx0.98$, \vs ${\approx}1.41$ for random orthogonal
maps), and our pre-registered comparison is \emph{negative}: augmentation does not
improve cyclicity ($3$ seeds per regime; $0.983$ \vs $0.981$; one-sided $p=0.83$), so
DR-H2b is not supported.

\begin{figure}[t]
\centering
\includegraphics[width=0.5\columnwidth]{fig_fifths_geometry.pdf}
\caption{The twelve major-tonic probe directions (L$4$, first two PCs) self-organize
into a circle of fifths. No key, chord, or degree token exists in the vocabulary.}
\label{fig:fifths}
\end{figure}

\subsection{Editing the key state changes the composed key}
\label{sec:h3}
At the best in-budget condition (V-PROBE, L$4$),
$\tkr(\text{edit})-\tkr(\text{K1})>0$ holds for \textbf{$12/12$} targets
(Holm-corrected $p<10^{-4}$; rank-biserial $r=0.64$--$1.00$) against a
pre-registered bar of ${\geq}\,8/12$: \textbf{DR-H3 is supported}
(Fig.~\ref{fig:bars}). Pitch content triangulates the same conclusion without the
estimator: $\ikr_{\kappa^{*}}{=}0.937$ \vs $\ikr_{\kappa_{\mathrm{src}}}{=}0.646$
after editing. Success falls gently with fifths distance ($0.395$ at distance~$1$ to
$0.280$ at the tritone), while random edits collapse to zero once the identity
target is excluded. The edit reaches $58\%$ of the behavioral ceiling (K4): the state
is readable in a rank-$24$ subspace of one layer, yet a single-layer edit realizes
only part of the behavioral effect---an \emph{actionability gap}. (Guarded edit over unguarded
K4, whose perplexity twin is untransposed and so scores transposition, not damage;
guarding both sides gives $62\%$, neither $66\%$---we report the smallest.) Two subsidiary
findings: the \emph{optimized} subspace (V-DAS, $0.213$) underperforms the probe's
readout directions, and V-MEAN edits destroy quality at deep layers (guard pass
${\leq}5\%$)---the guard is load-bearing, as Sec.~\ref{sec:capacity} makes decisive.

\begin{figure}[t]
\centering
\includegraphics[width=0.96\columnwidth]{fig_intervention_bars.pdf}
\caption{Guarded strict $\tkr$ by condition: the continuation lands exactly on the
injected key \emph{and} stays inside the frozen musicality budget. Chance $=1/12$.
Bars: prompt-level bootstrap $95\%$ CI.}
\label{fig:bars}
\end{figure}

\subsection{Readability is diffuse; causal efficacy is sharp}
The edit-minus-control effect is single-peaked across layers
($0.03,0.04,0.19,0.26,\mathbf{0.30},0.26,0.23,0.23$ for L$0$--L$7$); the
best-\vs-median bootstrap CI $[0.038,0.109]$ excludes zero (DR-H5 supported). The
two profiles in Fig.~\ref{fig:layers} are \emph{not} the same shape, and that is the
finding. The probe saturates across L$2$--L$6$ (macro-$F_1$ $0.920$--$0.929$, all
within $0.01$ of maximum): the key is legible from most of the network. Yet a
single-layer edit moves the output only near L$4$, and barely at all at
L$0$--L$1$. \emph{Where a state is legible is a weak guide to where it is used}---a
caution for work that localizes a computation by probing accuracy alone.

\begin{figure}[t]
\centering
\includegraphics[width=0.92\columnwidth]{fig_layer_profile.pdf}
\caption{(a)~Probe macro-$F_1$ per layer, all six M-CTRL models (thin lines; R-Aug
seed~$0$ highlighted): the key is legible across a broad plateau. (b)~Causal effect:
guarded $\tkr$ of the V-PROBE edit \vs the K1 matched random control (bands:
prompt-level bootstrap $95\%$ CI). Only a narrow band around L$4$ moves the output.}
\label{fig:layers}
\end{figure}

\subsection{What a world model actually requires}
\label{sec:capacity}
Training the same data and protocol at four capacities (L$2$/$d128$, L$4$/$d256$,
L$8$/$d512$, L$12$/$d768$; $0.5$M, $3.4$M, $26$M, $86$M parameters counted from the
checkpoints) dissociates three things the term ``world model'' conflates
(Table~\ref{tab:what}, top). \emph{Behaviour} is flat: next-token top-$1$ is
$0.868$--$0.879$ across this $170\times$ range. \emph{Legibility} is not: at $0.5$M the
probe peaks \emph{below} the input baseline and DR-H1 is unsupported in both seeds---a
behaviourally equivalent model with no key state, the surface-statistics account made
flesh, which the pre-registered test correctly refuses to credit. \emph{Manipulability}
is a third thing again: at $0.5$M the edit moves the key \emph{more} often than at $26$M
in raw terms ($0.55$ \vs $0.43$) and survives the musicality guard only $10\%$ of the
time (\vs $80\%$)---with two layers, ``editing the key subspace'' is a blunt overwrite,
not the movement of a separable state. Reporting guarded $\tkr$ alone hides this;
reporting raw $\tkr$ alone credits a surface model with a world model. The guard is what
tells them apart. Nor does manipulability simply keep improving: it peaks near $26$M and
falls to $0.18$ at $86$M, and a full twelve-layer sweep confirms this is not a
layer-choice artifact (the $86$M causal peak, L$6$ at $0.206$, is indistinguishable from
the L$3$ we had used, $0.202$).

\subsection{A public model trained on real music: it reads, and it uses}
\label{sec:wild}
Everything above is our own model on our own data. We close by leaving both. On $300$
Bach chorales whose \emph{local} key at every note comes from published Roman-numeral
analyses,\footnote{Scores: \texttt{bach-370-chorales} (CC BY-NC-SA); analyses:
\emph{When in Rome} (CC BY-SA). Neither is redistributed.} our $86$M model's key
subspace transfers ($F_1$ $0.534$ \vs a $0.035$ selectivity floor) but does \emph{not}
beat the pitch surface: the corrected margin, $+0.034$ (CI $[-0.021,0.103]$), does not
exclude zero.

The Anticipatory Music Transformer~\cite{thickstun2023anticipatory} (Apache-2.0; Lakh
MIDI, MetaMIDI, FMA, $450$k commercial records) is, at its \texttt{small} size,
\emph{architecturally identical} to that model---$12$ layers, $d{=}768$, $85.1$M
\vs $85.2$M non-embedding parameters (its $128$M total only reflects a $55$k
vocabulary against our $124$)---and its (arrival-time, duration, note) vocabulary is
leak-free exactly as ours is.
Probed on the same chorales, at the same positions, with the same controls and rule (the
two C3 baselines land at $0.451$ and $0.454$, confirming the setups match), it reaches
$F_1$ $0.690$ with a corrected margin of $+0.196$ (CI $[0.132,0.260]$): \textbf{DR-H1
supported}. With architecture, corpus, protocol and rule fixed, \emph{the training data
is the single free variable}, and it decides whether a key state is worth more than the
notes themselves.

\textbf{And that state is used.} We ran the same intervention on it: the probe's row
space as the subspace, a \emph{different} public checkpoint as the guard's reference
(medium refereeing small---different weights, not circular; $\ppl{=}0.849$ nats frozen
from $1{,}180$ natural modulations in the chorales), the edit layer searched on $20$
chorale prefixes and judged on $60$ it never saw, sham edit bit-identical ($20/20$).
\textbf{DR-H3 is supported}: at L$8$ the edit lands the continuation on the injected key
at $5.7\times$ the matched random control ($0.365$ \vs $0.064$), significant in $11/12$
targets, with $\ikr$ crossing ($0.824$ to the injected key \vs $0.785$ to the prompt's)
and---unlike our own model's $87\%$---\textbf{$100\%$ of edits inside the frozen
musicality budget}. A key state learned from real music is \emph{more} cleanly separable,
not less.

Two things only a real model could show. The causal peak (L$8$) is again not the probe
peak (L$10$), so Sec.~\ref{sec:h3}'s dissociation is not an artifact of synthetic data.
And the model \emph{resists}: the one target it cannot reach is F$\sharp$ major ($\tkr$
$0.000$), and across all twelve targets success is ordered almost perfectly by how common
that tonic is in the chorale corpus, counted over its note-weighted local-key labels
(Spearman $+0.91$, $n{=}12$, one-sided permutation $p<10^{-4}$; $+0.75$ if keys are
counted once per chorale instead). A model trained on real music can be steered into the
keys that music favours and resists the rest---a prior our uniformly-keyed synthetic
corpus could not have revealed. The chorale distribution is a proxy for tonal key
frequency, not a measurement of this model's own training corpus (Lakh MIDI and
MetaMIDI, not chorales), so the prior's \emph{source} is not identified here---only its
existence. One caution earned
the hard way: the probe must be read where the model is about to \emph{choose} a pitch.
Read at the note instead, the same model on the same chorales decays with depth
($F_1$ $0.545$ at L$0$ to $0.093$ at L$10$) rather than peaking deep, and the margin
falls to $+0.042$ (CI $[-0.034,0.092]$): DR-H1 would be \emph{unsupported}. One position
choice, and the verdict inverts.

\begin{table}[t]
\setlength{\tabcolsep}{3.5pt}\renewcommand{\arraystretch}{1.12}
\centering\small
\caption{What a world model requires. \textbf{Top}: same data and protocol, four
capacities --- behaviour is flat while legibility and manipulability are not.
\textbf{Bottom}: same architecture ($12$ layers, $d{=}768$), same chorales, same
protocol; only the training data differs. The public model both reads the key and
uses it.}
\label{tab:what}
\begin{tabularx}{\linewidth}{@{}Xcccc@{}}
\toprule
\textbf{Model} & \textbf{top-1} & \textbf{probe} & \textbf{DR-H1} & \textbf{guard}\\
\midrule
\multicolumn{5}{@{}l}{\emph{capacity (D-SYN, our models)}}\\
$0.5$M & $0.870$ & below C3 & no & $10\%$\\
$3.4$M & $0.878$ & $+0.05$ & yes & $60\%$\\
$26$M & $0.879$ & $+0.10$ & yes & $80\%$\\
$86$M & $0.879$ & $+0.10$ & yes & $71\%$\\
\midrule
\multicolumn{5}{@{}l}{\emph{real music, matched architecture (Bach chorales)}}\\
ours, synthetic-trained & --- & $+0.034$ & \textbf{no} & ---\\
public, real-trained~\cite{thickstun2023anticipatory} & --- & $\mathbf{+0.196}$ & \textbf{yes} & ---\\
\bottomrule
\end{tabularx}
\end{table}

\section{Discussion and Limitations}
\label{sec:disc}
The evidence supports the world-model reading within this training condition: a
Transformer trained on notes alone maintains a mid-layer key variable with
music-theoretic geometry, and generation consults it. The gap to the behavioral ceiling
and the failure of edits at L$0$--L$1$ point to a distributed, redundantly re-estimated
state rather than a single canonical register---a contrast with Othello-GPT's
near-complete editability that we attribute to music's continual re-grounding of
tonality in the preceding notes. K3 agrees: it is a negative control only at L$4$, where
it borrows L$0$'s failed subspace and does nothing ($0.064$, at or below K1), but
subspaces borrowed from L$5$--L$7$ and injected at L$1$--L$3$ reach $0.27$--$0.39$, so
the subspace is largely \emph{shared} across the middle stack and its layer-pooled mean
($0.160$) is not a ``partial effect''. The full version resolves K3 by layer.

Bounds. \emph{Mode:} the pre-registration fixed the injected mode to major; minor
targets were designated secondary and have not been run, so every causal claim here
concerns major targets. \emph{V-DAS below V-PROBE:} the optimized subspace may be
undertrained relative to its search space; we report it as observed. \emph{Estimator:}
the KS judge confuses fifths-neighbours, which strict/tolerant reporting brackets and
$\ikr$ triangulates. \emph{Capacity:} the sweep varies depth and width together,
locating a threshold, not its cause. \emph{Selection:} the best (method, layer) is
chosen on the data the Wilcoxon then tests, as pre-registered. \emph{Real models:} the
causal result holds for one public checkpoint at one size. The journal version carries
the ambiguity strata, the specificity matrix, the fifths-distance curve, per-seed
figures, and whether edited continuations modulate \emph{musically}. Training data are
synthetic; evaluation corpora are used under their stated licenses.

\section{Conclusion}
A music Transformer tracks musical key as an internal state variable that generation
consults---in our models, and in one people actually use. Capacity supplies not the key
but a key that can be moved without breaking the music; the training data, not the
architecture, decides whether that state beats the notes themselves. The testbed is
released as a benchmark.

\vspace{2pt}\noindent\footnotesize\textbf{Reproducibility.} Anonymized code, data
generator, pre-registration, and an audio demo page (clean \vs edited continuations):
\url{https://anonymous.4open.science/r/TBD}. \TBD{create before submission.}\normalsize

\begin{thebibliography}{9}\footnotesize
\bibitem{li2023othello} K.~Li, A.~K. Hopkins, D.~Bau, F.~Vi\'egas, H.~Pfister, and
M.~Wattenberg, ``Emergent world representations: Exploring a sequence model trained on
a synthetic task,'' in \emph{ICLR}, 2023.
\bibitem{nanda2023emergent} N.~Nanda, A.~Lee, and M.~Wattenberg, ``Emergent linear
representations in world models of self-supervised sequence models,'' in
\emph{BlackboxNLP}, 2023.
\bibitem{karvonen2024chess} A.~Karvonen, ``Emergent world models and latent variable
estimation in chess-playing language models,'' in \emph{CoLM}, 2024.
\bibitem{maze2024} A.~F. Spies, W.~Edwards, M.~I. Ivanitskiy, et al., ``Transformers use
causal world models in maze-solving tasks,'' \emph{arXiv:2412.11867}, 2024.
\bibitem{iris2026} X.~Zhang, ``What do world models learn in RL? Probing latent
representations in learned environment simulators,'' \emph{arXiv:2603.21546}, 2026.
\bibitem{syntheory2024} M.~Wei, M.~Freeman, C.~Donahue, and C.~Sun, ``Do music generation
models encode music theory?,'' in \emph{ISMIR}, 2024.
\bibitem{smitin2024} I.~Prokopiou, P.~Vikatos, M.~Kaliakatsos-Papakostas, T.~Giannakopoulos,
and T.~Stafylakis, ``Latent space disentanglement via activation steering for
interpretable attribute control in symbolic music generation,'' in \emph{EUSIPCO}, 2026.
\bibitem{thickstun2023anticipatory} J.~Thickstun, D.~Hall, C.~Donahue, and P.~Liang,
``Anticipatory music transformer,'' \emph{TMLR}, 2024.
\bibitem{krumhansl1982} C.~L. Krumhansl and E.~J. Kessler, ``Tracing the dynamic changes
in the perceived tonal organization in a spatial representation of musical keys,''
\emph{Psychological Review}, vol.~89, no.~4, pp.~334--368, 1982.
\bibitem{hewitt2019control} J.~Hewitt and P.~Liang, ``Designing and interpreting probes
with control tasks,'' in \emph{EMNLP}, 2019.
\end{thebibliography}

\end{document}
